% Ícones e diagramas vetoriais para o encerramento, sem telas fictícias de EDA.
\newcommand{\digitalTrilha}{%
\begin{circuitikz}[font=\scriptsize]
\node[nand port,scale=.7] (n) at (0,0) {};
\draw (n.in 1)--++(-.35,0);\draw (n.in 2)--++(-.35,0);
\draw (n.out)--(1.2,0);
\draw[blue_theme,fill=blue_theme!7] (1.2,-.65) rectangle (2.2,.65);
\node at (1.7,0){Reg.};\draw[seta] (2.2,0)--(2.8,0);
\draw[blue_theme,fill=ativo!8] (2.8,-.65) rectangle (3.8,.65);\node at (3.3,0){ULA};
\path[use as bounding box] (-1.2,-1) rectangle (4.2,1);
\end{circuitikz}}
\newcommand{\arquiteturaTrilha}{%
\begin{tikzpicture}[font=\scriptsize]
\draw[blue_theme,fill=blue_theme!7] (0,-.8) rectangle (1.7,.8);
\foreach \x in {.2,.9}{\foreach \y in {-.5,.15}{\fill[blue_theme!45] (\x,\y) rectangle ++(.5,.4);}}
\node[above] at (.85,.8){Processador};
\draw[<->,blue_theme,thick] (1.7,0)--(3,0);
\foreach \y in {-.7,-.2,.3}{\draw[blue_theme,fill=ativo!10] (3,\y) rectangle ++(1.6,.35);}
\node[above] at (3.8,.8){Memória};
\path[use as bounding box] (-.4,-1.2) rectangle (5,1.3);
\end{tikzpicture}}
\newcommand{\analogicaTrilha}{%
\begin{circuitikz}[font=\scriptsize]
\node[nmos] (n) at (0,0) {};
\draw (n.D)--(0,.9) to[R,l=$R_D$] (0,2.6) node[above]{$\VDD$};
\draw (n.S)--(0,-.8) node[ground]{};
\draw (n.G)--++(-.4,0) node[left,align=center]{$V_{\mathrm{bias}}$\\$+v_{\mathrm{in}}$};
\draw (0,.9) to[short,*-] (1,.9) node[right]{$v_{\mathrm{out}}$};
\path[use as bounding box] (-2,-1.1) rectangle (1.8,3.1);
\end{circuitikz}}
\newcommand{\mistoTrilha}{%
\begin{tikzpicture}[font=\scriptsize,x=.8cm,y=.8cm]
\draw[blue_theme,thick,domain=0:4,samples=65] plot (\x,{1+.7*sin(90*\x)});
\draw[ativo,thick] (0,1)--(.5,1)--(.5,1.5)--(1,1.5)--(1.5,1.5)--(2,1.5)--(2,1)--(2.5,1)--(2.5,.5)--(3,.5)--(3.5,.5)--(4,.5);
\foreach \x in {0,.5,...,3.5}{\fill[ativo] (\x,{.5*round(2*(1+.7*sin(90*\x)))}) circle(.045);}
\node[blue_theme] at (2,2.4){Contínuo};\node[ativo] at (2,-.35){Amostras e códigos};
\path[use as bounding box] (-.2,-.8) rectangle (4.3,2.9);
\end{tikzpicture}}
\newcommand{\rfTrilha}{%
\begin{circuitikz}[font=\scriptsize]
\draw (0,1) to[L,l=$L$] (1.7,1)--(2.5,1);
\draw (1.7,1) to[C,l=$C$] (1.7,-.6) node[ground]{};
\draw[blue_theme,thick] (2.5,1)--(2.5,2)--(2.05,2.4) (2.5,2)--(2.95,2.4);
\draw (0,1)--(-.3,1) node[left]{RF};
\path[use as bounding box] (-.8,-1) rectangle (3.3,2.8);
\end{circuitikz}}
\newcommand{\produtoIntegrado}{%
\begin{tikzpicture}[x=1cm,y=1cm,font=\scriptsize]
% Uma planta funcional, sem alegar geometria de um produto real.
\draw[blue_theme,thick,fill=blue_theme!3] (0,0) rectangle (10,4.3);
\foreach \x in {.5,1.3,...,9.5}{\draw[blue_theme] (\x,0)--(\x,-.25) (\x,4.3)--(\x,4.55);}
\node[blue_theme,font=\small\bfseries] at (5,5){Um sistema integrado};
\fill[blue_theme!14] (.3,.4) rectangle (2.5,3.8);\node[align=center,text width=2cm] at (1.4,3.2){\textbf{Analógica}\\Sensor e condicionamento};
\draw[blue_theme,thick,domain=.55:2.2,samples=40] plot (\x,{1.8+.4*sin(360*\x)});
\fill[lacuna!18] (2.8,.4) rectangle (4.4,3.8);\node[align=center] at (3.6,3.2){\textbf{Sinais mistos}\\ADC};
\draw[lacuna,thick] (3,1.5)--(3.3,1.5)--(3.3,1.8)--(3.6,1.8)--(3.6,2.1)--(3.9,2.1)--(3.9,2.4)--(4.2,2.4);
\fill[ativo!13] (4.7,.4) rectangle (7.3,3.8);\node[align=center,text width=2.35cm] at (6,3.2){\textbf{Digital e arquitetura}\\Processamento e memória};
\foreach \x in {5,5.65,6.3}{\foreach \y in {1.3,2}{\draw[ativo] (\x,\y) rectangle ++(.45,.4);}}
\fill[blue_theme!10] (7.6,.4) rectangle (9.7,3.8);\node[align=center] at (8.65,3.2){\textbf{RF}\\Comunicação};
\draw[blue_theme,thick] (8.6,1.2)--(8.6,2.2)--(8.2,2.6) (8.6,2.2)--(9,2.6);
\foreach \x/\z in {2.5/2.8,4.4/4.7,7.3/7.6}{\draw[seta] (\x,2.7)--(\z,2.7);}
\node[font=\scriptsize,black!65] at (5,-.75){Layout · verificação · processos · teste · encapsulamento · ferramentas EDA};
\path[use as bounding box] (-.4,-1.1) rectangle (10.4,5.4);
\end{tikzpicture}}
